\subsection{三点固定部分の幾何学的monodromy}
\label{subsec:z1-fixed-monodromy-v26}

\begin{proposition}[固定三点への射影は全対称群]
\label{prop:z1-fixed-monodromy-v26}
$Z_1$ の被約二十一点配置上の幾何学的monodromyを
\[
 G_{Z_1}\subset
 \mathfrak S_3\times\bigl((\mathbb Z/2\mathbb Z)^9\rtimes\mathfrak S_9\bigr)
\]
と書くと、第一因子への射影は $\mathfrak S_3$ 全体である。
この命題は $\mathfrak S_3\times\{1\}\subset G_{Z_1}$ を主張しない。
\end{proposition}

\begin{proof}
ここで monodromy の基底は指定反対合を持つ枠付き成分であり、粗有理商ではない。
固定点を定める $B$ の退化スキームは、基点 $B_-$ で既に有限被約である。
正固有行を $u_{12},u_{13},u_{23}$、負固有行を $u_{11},u_{22},u_{33}$ とする順序では、
基点の負ブロックは
\[
 B_-(x,y,z)=
 \begin{pmatrix}
 0&0&0&y&0&z&0\\
 0&x&0&0&z&0&0\\
 x&0&y&0&0&0&0
 \end{pmatrix}.
\]
三行の台となる列集合は互いに交わらないので、退化点は三つの座標点だけである。
例えば $x\ne0$ の座標点では、二つの三次小行列式が $x^2y$ と $x^2z$ なので、
局所イデアルは $(y,z)$ となる。他の二点も同様である。
実際、三次小行列式のイデアルは
$(x,y,z)(xy,xz,yz)$ で、その飽和は $(xy,xz,yz)$ である。

従って $Z_1$ の滑らかな開部分 $W$ を、$B_-$ を含み、
$B$ の退化スキーム $\mathcal T\to W$ が有限エタールな三重被覆となるように取れる。
この延長の平坦性も確認できる。固有性から基点近傍で有限となり、
一般階数は三、被約な特殊ファイバーの次元も三である。
正則局所整域上で Nakayama の補題による $R^3$ から有限代数への全射を取ると、
核は自由加群内の階数零の部分加群なので零である。
従って有限自由階数三であり、被約な特殊ファイバーから近傍でエタールとなる。
以下の枠変更群の基点軌道も $W$ に含めてよい。
その軌道では退化スキームは上の三点の射影変換であり、
滑らかさと有限エタール性が保たれるためである。

$J_0=\operatorname{diag}(-1,1,1,-1,1,-1)$ を元の中央矢順序で書く。
同じ平面座標変更を左右に施し、中央矢の基底変更を $J_0$ の中心化群に取ると、
指定反対合を保つ連結群
\[
 \mathrm{GL}_3\times
 \bigl(\mathrm{GL}_3\times\mathrm{GL}_3\bigr)
\]
が作用する。関係枠を保つ表示では、さらに連結群 $\mathrm{GL}_7$ を付けてもよい。
この群は $Z_1$ を保つ。

三巡回置換 $\sigma=(123)$ に対する
$(\sigma,\operatorname{Sym}^2\sigma)$ と、互換 $\rho=(12)$ に対する
$(\rho,J_0\operatorname{Sym}^2\rho)$ は、ともに $B_-$ を固定する。
これを左関係の列 $r_0,\ldots,r_6$ で確認すると、像はそれぞれ
\[
 (r_3,r_2,r_5,r_4,r_0,r_1,r_6),\qquad
 (-r_2,-r_3,-r_0,-r_1,-r_5,-r_4,r_6)
\]
となる。中央作用はいずれも $J_0$ と可換であり、右関係も同時に保たれる。
関係pairingについても、対応する七次行列 $M$ は $MAM^t=A$ を満たす。

複素数上でこの連結枠変更群の恒等元から各安定化元まで経路を取り、
$B_-$ に作用させると、$W$ 内の基点軌道上の閉路を得る。
その三点被覆への作用は座標点の三巡回置換と互換である。
従って $\mathcal T\to W$ の幾何学的monodromyは $\mathfrak S_3$ である。

最後に、全二十一点が有限エタールになる非空開部分を $U\subset W$ とする。
滑らかな既約基底の稠密開部分への制限で、有限エタール被覆のmonodromyは縮まらない。
例えば被覆の連結Galois閉包は既約であり、その非空開部分も連結なことから従う。
よって $\mathcal T|_U$ のmonodromyも $\mathfrak S_3$ である。
これは全二十一点のmonodromyの固定三点への射影に一致する。
\end{proof}

\paragraph{再現入口.}
\nolinkurl{z1_fixed_monodromy_v26.py} と
\nolinkurl{z1_fixed_monodromy_v26_certificate.json} が、
三次小行列式の単項式生成元、二つの中央行列と関係行列、
$J_0$ との可換性、関係およびpairingの保存を有理数上で厳密に確認する。
幾何学的monodromyへの移行は、上の連結群軌道と稠密開制限の証明による。
